\documentclass[a4paper, UTF-8]{article}
\usepackage[margin=1in]{geometry}
\usepackage{algorithm}
\usepackage{algpseudocode}
\usepackage{amsmath}
\usepackage{amsfonts}
\usepackage{amsthm}

\usepackage[utf8]{inputenc} % allow utf-8 input
\usepackage[T1]{fontenc}    % use 8-bit T1 fonts
\usepackage{hyperref}       % hyperlinks
\usepackage{url}            % simple URL typesetting
\usepackage{booktabs}       % professional-quality tables
\usepackage{amsfonts}       % blackboard math symbols
\usepackage{nicefrac}       % compact symbols for 1/2, etc.
\usepackage{microtype}      % microtypography
\usepackage{xcolor}         % colors

\theoremstyle{plain}
\newtheorem{theorem}{\indent Theorem}[section]
\newtheorem{lemma}[theorem]{\indent Lemma}
\newtheorem{assumption}{\indent Assumption}[section]
\newtheorem{note}[theorem]{\indent Notation}
\newtheorem{proposition}[theorem]{\indent Proposition}
\newtheorem{corollary}[theorem]{\indent Corollary}
\newtheorem{definition}{\indent Definition}[section]
\newtheorem{example}{\indent Example}[section]
\newtheorem{remark}{\indent Remark}[section]
\newenvironment{solution}{\begin{proof}[\indent\bf Solution]}{\end{proof}}
\renewcommand{\proofname}{\indent\bf Proof}

\begin{document}

\title{\bf{Assignment 3}}

\author{\bf{Yanqiao Chen}\\ 
  Department of Computer Science and Engineering\\
  Southern University of Science and Technology\\
  \texttt{12412115@mail.sustech.edu.cn}
  }


\maketitle

\section{Page 30, Problem 46}

\subsection{a}

Let the event $A$ be $\{\text{a red ball is drawn}\}$, and the event $B$ be $\{\text{the coin shows heads up}\}$.

Then we have:
\begin{align*}
P(A) &= P(A\mid B)P(B) + P(A\mid \bar{B})P(\bar{B})\\
     &= \frac{3}{5}\cdot \frac{1}{2} + \frac{2}{7}\cdot \frac{1}{2} \\
     &= \frac{31}{70}
\end{align*}

\subsection{b}

From the Bayes' theorem, we have:
\begin{align*}
P(B|A) &= \frac{P(A|B)P(B)}{P(A)} \\
       &= \frac{\frac{3}{5}\cdot \frac{1}{2}}{\frac{31}{70}} \\
       &= \frac{21}{31}
\end{align*}

\section{Page 31, Problem 53}

Let the event: 

\begin{itemize}

\item $A_1$ be $\{\text{the claims come from the high risk client}\}$
\item $A_2$ be $\{\text{the claims come from the medium risk client}\}$
\item $A_3$ be $\{\text{the claims come from the low risk client}\}$
\item $B$ be $\{\text{the claims comes}\}$.
\end{itemize}

\begin{align*}
P(A_1\mid B) &= \frac{P(B\mid A_1)P(A_1)}{P(B\mid A_1)P(A_1) + P(B\mid A_2)P(A_2) + P(B\mid A_3)P(A_3)}\\
             &= \frac{0.02 \cdot 0.1}{0.02 \cdot 0.1 + 0.01 \cdot 0.2 + 0.0025 \cdot 0.7}\\
             &= \frac{0.002}{0.002 + 0.002 + 0.00175} \\
             &= \frac{8}{23}
\end{align*}


\section{Page 31, Problem 54}

\subsection{a}

\begin{align*} 
P(R_{i+1}) &= P(R_{i+1}\mid R_i)P(R_i) + P(R_{i+1}\mid \bar{R}_i)P(\bar{R}_i) \\
           &= \alpha \cdot p + (1-\beta) \cdot (1-p)\\
           &= (\alpha + \beta - 1) \cdot p + (1-\beta)
\end{align*}

\subsection{b} 

\begin{align*}
P(R_{i+2}) &= P(R_{i+2}\mid R_{i+1})P(R_{i+1}) + P(R_{i+2}\mid \bar{R}_{i+1})P(\bar{R}_{i+1}) \\
           &= \alpha \cdot P(R_{i+1}) + (1-\beta) \cdot (1-P(R_{i+1})) \\
           &= (\alpha + \beta - 1) \cdot (\alpha + \beta - 1) \cdot p + (1-\beta) \cdot (\alpha + \beta - 1) + (1-\beta)\\
           &= (\alpha + \beta - 1)^2 \cdot p + (1-\beta) \cdot (\alpha + \beta)
\end{align*}

\subsection{c} 

This is a linear homogeneous recurrence relation $p_{i + 1} = (\alpha + \beta - 1)p_i + (1-\beta)$. 

Using undetermined coefficients, 
$$
(p_{i + 1} + k) = (\alpha + \beta - 1)(p_i + k) \Leftrightarrow (\alpha + \beta - 2)k = 1-\beta
$$

Thus, we have $k = \frac{1-\beta}{\alpha + \beta - 2}$.

Then we have:
$$
p_i + k = (\alpha + \beta - 1)^i (p + k)
$$

Thus, we conclude that:
$$
p_n = (\alpha + \beta - 1)^n (p + \frac{1-\beta}{\alpha + \beta - 2}) - \frac{1-\beta}{\alpha + \beta - 2}
$$

From the property of the probability, $P(R_{i}|R_{i-1}) + P(\bar{R}_{i}|\bar{R_{i-1}}) = \alpha + \beta < 2$, then $\alpha + \beta - 1 < 1$. According to the property of the geometric series, the $\lim_{n\to \infty} p_n$ converges to $\frac{1-\beta}{2 - \alpha - \beta}$.


\section{Page 32, Problem 63}

Let 
\begin{itemize}
  \item event $A$ be $\{\text{a person reaches her 70th birthday}\}$
  \item event $B$ be $\{\text{a person reaches her 80th birthday}\}$.
\end{itemize}

Then we have:
\begin{align*}
P(B\mid A) = \frac{P(AB)}{P(A)} 
\end{align*}

As $B \subseteq A$, we have $P(AB) = P(B)$, so:
$$
P(B\mid A) = \frac{P(B)}{P(A)} = \frac{0.2}{0.6} = \frac{1}{3}
$$

\section{Supplementary Problem 1.1}

The Monty Hall problem is a famous problem in probability theory. The problem is as follows:

\begin{quote}
You are a contestant on a game show, and you are given the choice of three doors. Behind one door is a car; behind the others, goats. You pick a door, say No. 1, and the host, who knows what's behind the doors, opens another door, say No. 3, which has a goat. He then says to you, "Do you want to switch to door No. 2?" Is it to your advantage to switch your choice?
\end{quote}

Let

\begin{itemize}
  \item event $A_i$ be $\{\text{the contestant choose door No. } i\}$
  \item event $B_i$ be $\{\text{the host opens door No. } i\}$.
  \item event $C_i$ be $\{\text{the car is behind door No.} i\}$.
\end{itemize}

We want to find $P(C_2\mid A_1B_3)$.

\begin{align*} 
P(C_2\mid A_1B_3) &= \frac{P(A_1B_3\mid C_2)P(C_2)}{P(A_1B_3\mid C_1)P(C_1) + P(A_1B_3\mid C_2)P(C_2) + P(A_1B_3\mid C_3)P(C_3)} \\
&= \frac{1 \cdot \frac{1}{3}}{\frac{1}{2} \cdot \frac{1}{3} + 1 \cdot \frac{1}{3} + 0 \cdot \frac{1}{3}} \\
&= \frac{\frac{1}{3}}{\frac{1}{6} + \frac{1}{3}} \\
&= \frac{\frac{1}{3}}{\frac{1}{2}} \\
&= \frac{2}{3}
\end{align*}

Then the contestant should switch to door No. 2, which has a higher probability to win the car. $\square$

\section{Supplementary Problem 1.2}

Let 
\begin{itemize}
  \item event $A$ be $\{\text{the child get sick}\}$.
  \item event $B$ be $\{\text{the mother get sick}\}$.
  \item event $C$ be $\{\text{the father get sick}\}$.
\end{itemize}

We have 

\begin{itemize}
  \item $P(A) = 0.6$
  \item $P(B \mid A) = 0.5$
  \item $P(C \mid AB) = 0.4$
\end{itemize}

We would like to find the probability 
$$
P(AB\bar{C}) = P(\bar{C}\mid AB)P(AB) = (1 - P(C\mid AB))P(B\mid A)P(A) = (1 - 0.4) \cdot 0.5 \cdot 0.6 = 0.18
$$

\section{Supplementary Problem 1.3}

Let 

\begin{itemize}
  \item event $A$ be $\{\text{the first product of the passed inspection}\}$.
  \item event $B$ be $\{\text{the machine is well-adjusted}\}$.
\end{itemize}

We compute the 
\begin{align*}
P(B|A) &= \frac{P(A|B)P(B)}{P(A)} \\
       &= \frac{P(A|B)P(B)}{P(A|B)P(B) + P(A|\bar{B})P(\bar{B})} \\
       &= \frac{0.98\cdot 0.95}{0.98\cdot 0.95 + 0.05\cdot 0.55}\\
       &= 0.9713
\end{align*}

\section{Page 32, Problem 68}

The statement is false. Independence is not transitive. We provide a counterexample.

Consider the sample space $\Omega = \{1, 2, 3, 4, 5, 6\}$ with uniform probability distribution.

Define three events:
\begin{itemize}
    \item $A = \{1, 2, 3, 4\}$, so $P(A) = \frac{4}{6} = \frac{2}{3}$
    \item $B = \{1, 2, 5\}$, so $P(B) = \frac{3}{6} = \frac{1}{2}$
    \item $C = \{1, 2, 3, 6\}$, so $P(C) = \frac{4}{6} = \frac{2}{3}$
\end{itemize}

\textbf{Verification:}

1. Check $A$ and $B$:
   - $A \cap B = \{1, 2\}$, so $P(AB) = \frac{2}{6} = \frac{1}{3}$
   - $P(A)P(B) = \frac{2}{3} \times \frac{1}{2} = \frac{1}{3}$
   - Therefore, $A$ and $B$ are independent.

2. Check $B$ and $C$:
   - $B \cap C = \{1, 2\}$, so $P(BC) = \frac{2}{6} = \frac{1}{3}$
   - $P(B)P(C) = \frac{1}{2} \times \frac{2}{3} = \frac{1}{3}$
   - Therefore, $B$ and $C$ are independent.

3. Check $A$ and $C$:
   - $A \cap C = \{1, 2, 3\}$, so $P(AC) = \frac{3}{6} = \frac{1}{2}$
   - $P(A)P(C) = \frac{2}{3} \times \frac{2}{3} = \frac{4}{9}$
   - Since $\frac{1}{2} \neq \frac{4}{9}$, $A$ and $C$ are \textbf{not} independent.

Thus, we have $A \perp B$ and $B \perp C$, but $A \not\perp C$, proving that independence is not transitive.

\section{Page 32, Problem 71}

From the property of mutually independent events, we have:
\begin{align*}
P(ABC) &= P(A)P(B)P(C) \\
P(AB) &= P(A)P(B) \\
P(AC) &= P(A)P(C) \\
P(BC) &= P(B)P(C)
\end{align*}

Thus, 
\begin{align*}
P(ABC) = P(A)P(B)P(C) = P(AB)P(C)
\end{align*}

Then we conclude that $P(ABC) = P(AB)P(C)$, which means that $AB$ is independent of $C$. Similarly, 

\begin{align*} 
P((A \cup B)C) &= P(AC \cup BC) = P(AC) + P(BC) - P(ABC)\\ 
               &= P(A)P(C) + P(B)P(C) - P(A)P(B)P(C)\\ 
               &= (P(A)+P(B)-P(AB))P(C) \\
               &= P((A \cup B))P(C)\\
\end{align*}

which means that $A \cup B$ is independent of $C$. $\square$

\section{Page 33, Problem 74}

Let $P(A_i)$ represent "the unit $i$ fails", with $P(A_i) = p$.

The system works if at least one of the three parallel branches works:
- Branch 1: units 1 and 2 in series (works if both work)
- Branch 2: unit 3 alone
- Branch 3: units 4 and 5 in series

Let $S_1 = \overline{A_1 \cup A_2}$ (branch 1 works), $S_2 = \overline{A_3}$ (branch 2 works), $S_3 = \overline{A_4 \cup A_5}$ (branch 3 works).

Since the branches share no common units, they are independent. We use the complement rule:

\begin{align*}
P(\text{system works}) &= 1 - P(\text{all branches fail}) \\
&= 1 - P(S_1^c)P(S_2^c)P(S_3^c)
\end{align*}

where:
\begin{align*}
P(S_1^c) &= P(A_1 \cup A_2) = 1 - P(\bar{A_1})P(\bar{A_2}) = 1 - (1-p)^2 = 2p - p^2 \\
P(S_2^c) &= P(A_3) = p \\
P(S_3^c) &= P(A_4 \cup A_5) = 1 - (1-p)^2 = 2p - p^2
\end{align*}

Therefore:
\begin{align*}
P(\text{system works}) &= 1 - (2p-p^2) \cdot p \cdot (2p-p^2) \\
&= 1 - p(2p-p^2)^2 \\
&= 1 - p(4p^2 - 4p^3 + p^4) \\
&= 1 - 4p^3 + 4p^4 - p^5
\end{align*}

Alternatively, using inclusion-exclusion principle:
\begin{align*}
P(S_1 \cup S_2 \cup S_3) &= P(S_1) + P(S_2) + P(S_3) - P(S_1S_2) - P(S_1S_3) - P(S_2S_3) + P(S_1S_2S_3) \\
&= 2(1-p)^2 + (1-p) - 2(1-p)^3 - (1-p)^4 + (1-p)^5 \\
&= 1 - 4p^3 + 4p^4 - p^5
\end{align*}

\section{Page 33, Problem 77}

Let the throwing time be $n$, let $A_i$ be "the player succeed".

We obtain 
\begin{align*}
P(\bigcup_{i=1}^n A_i)  &= 1- P(\bigcap_{i=1}^n \bar{A_i}) \\
                        &= 1 - (0.95)^n \\ 
                        &\geq 0.5
\end{align*}

Thus the inequality will be 
$$
(0.95)^n \leq 0.5
$$

Taking logarithm we obtain 
$$
n \geq \frac{\ln 0.5}{\ln(0.95)} \approx 13.51
$$

As $n \in \mathbb{Z}$, thus $n_{min} = 14$.

\section{Page 34, Problem 79}

\subsection{a} 

It is obvious that 
\begin{align*} 
P(AA) = \frac{1}{4} \\
P(Aa) = \frac{1}{2} \\ 
P(aa) = \frac{1}{4} 
\end{align*} 

\subsection{b} 

Let $B$ be "the male offspring is not diseased", $A$ be "the male is a carrier".

The probability is 
$$
P(A|B) = \frac{P(AB)}{P(B)} = \frac{2}{4} \times \frac{4}{3} = \frac{2}{3}
$$

\subsection{c} 

Let $F$ denote the father (the nondiseased offspring from part (b)), and $M$ denote the mother (a random population member with probability $p$ of being a carrier).

From part (b), we have:
\begin{align*}
P(F = AA \mid \text{nondiseased}) &= \frac{1}{3} \\
P(F = Aa \mid \text{nondiseased}) &= \frac{2}{3}
\end{align*}

For the mother:
\begin{align*}
P(M = Aa) &= p \\
P(M = AA) &= 1-p
\end{align*}

We calculate the probabilities by considering all four mating combinations:

\textbf{Case 1:} $F = AA$ and $M = Aa$ (probability $\frac{1}{3} \cdot p$)
\begin{itemize}
    \item Offspring $AA$: $\frac{1}{2}$
    \item Offspring $Aa$: $\frac{1}{2}$
    \item Offspring $aa$: $0$
\end{itemize}

\textbf{Case 2:} $F = Aa$ and $M = Aa$ (probability $\frac{2}{3} \cdot p$)
\begin{itemize}
    \item Offspring $AA$: $\frac{1}{4}$
    \item Offspring $Aa$: $\frac{1}{2}$
    \item Offspring $aa$: $\frac{1}{4}$
\end{itemize}

\textbf{Case 3:} $F = AA$ and $M = AA$ (probability $\frac{1}{3} \cdot (1-p)$)
\begin{itemize}
    \item Offspring $AA$: $1$
    \item Offspring $Aa$: $0$
    \item Offspring $aa$: $0$
\end{itemize}

\textbf{Case 4:} $F = Aa$ and $M = AA$ (probability $\frac{2}{3} \cdot (1-p)$)
\begin{itemize}
    \item Offspring $AA$: $\frac{1}{2}$
    \item Offspring $Aa$: $\frac{1}{2}$
    \item Offspring $aa$: $0$
\end{itemize}

Now we compute the total probabilities:

\begin{align*}
P(\text{offspring is } AA) &= \frac{1}{3}p \cdot \frac{1}{2} + \frac{2}{3}p \cdot \frac{1}{4} + \frac{1}{3}(1-p) \cdot 1 + \frac{2}{3}(1-p) \cdot \frac{1}{2}\\
&= \frac{p}{6} + \frac{p}{6} + \frac{1-p}{3} + \frac{1-p}{3}\\
&= \frac{p}{3} + \frac{2(1-p)}{3}\\
&= \frac{2-p}{3}
\end{align*}

\begin{align*}
P(\text{offspring is } Aa) &= \frac{1}{3}p \cdot \frac{1}{2} + \frac{2}{3}p \cdot \frac{1}{2} + \frac{1}{3}(1-p) \cdot 0 + \frac{2}{3}(1-p) \cdot \frac{1}{2}\\
&= \frac{p}{6} + \frac{p}{3} + 0 + \frac{1-p}{3}\\
&= \frac{p}{6} + \frac{2p}{6} + \frac{2(1-p)}{6}\\
&= \frac{p + 2}{6}
\end{align*}

\begin{align*}
P(\text{offspring is } aa) &= \frac{1}{3}p \cdot 0 + \frac{2}{3}p \cdot \frac{1}{4} + \frac{1}{3}(1-p) \cdot 0 + \frac{2}{3}(1-p) \cdot 0\\
&= 0 + \frac{p}{6} + 0 + 0\\
&= \frac{p}{6}
\end{align*}

We can verify that these probabilities sum to 1:
\[
\frac{2-p}{3} + \frac{p+2}{6} + \frac{p}{6} = \frac{4-2p+p+2+p}{6} = \frac{6}{6} = 1
\]

\section{Supplementary Problem 2.1}

Given that 
$$
P(\overline{A \cup B})= P(\bar{A})P(\bar{B})= \frac{1}{9}
$$
and 
$$
P(A \bar{B}) = P(\bar{A}B) \Leftrightarrow P(A)P(\bar{B}) = P(\bar{A}) P(B) = P(\bar{A})(1 - P(\bar{B})) = P(\bar{A}) - \frac{1}{9}
$$

Notice that 
$$
P(A \cup B) = P(A) + P(B) - P(A)P(B)
$$

Thus 
$$
P(A)P(\bar{B}) = P(B) - P(A)P(B)
$$

Then 
$$
P(A) = P(B)
$$

Thus 
$$
P(\bar{A}) = \frac{1}{3}
$$

And we conclude that 
$$
P(A) = \frac{2}{3}
$$

\section{Supplementary Problem 2.2}

\begin{align*} 
P(A\cup B\cup C) = P(A) + P(B) + P(C) - P(AB) - P(BC) - P(AC) = \frac{9}{16}
\end{align*}

Denote 
$$
P(A) = p
$$

We obtain 
$$
16p^2 - 16 p  +3 = 0
$$

Solving this 
$$
(4p - 3)(4p - 1) = 0
$$

We obatin 
$$
p_1 = \frac{3}{4}, p_2 = \frac{1}{4}
$$

However, as 
$$
P(AB) = P(AB \cap \bar{C}) \leq P(\bar{C}) = 1-p
$$

Thus, 
$$
p^2 \leq 1 - p
$$

We obtain 
$$
p \leq \frac{\sqrt{5} - 1}{2} \approx 0.618 < \frac{3}{4}
$$

Thus 
$$
p = \frac{1}{4}
$$

\end{document}
